\subsection{Synthetic Task Generation}\label{sec:methodology-synth}

Real (image, mask) pairs are limited to the anatomical classes and
appearance distribution covered by the training cohort. To broaden
training beyond this, a fraction $p_{\text{synth}}$ of training tasks are
drawn from a synthetic generator instead of real context/target subjects.

\paragraph{Supervoxel repainting.} A real-anatomy volume is first
oversegmented into supervoxels (unsupervised, appearance-agnostic
regions). One supervoxel is treated as the target class, and its binary
mask is label-perfect by construction. Each of the $K{+}1$ volumes needed
for a task is then produced by repainting the same supervoxel geometry
with independently sampled per-class Gaussian intensities (a Gaussian
mixture over classes), giving $K{+}1$ (image, mask) pairs that share
identical anatomy but differ in appearance -- no real cross-subject
registration or annotation is required.

\paragraph{Paint noise.} The default painting noise is i.i.d.\ per voxel:
it matches each class's target intensity variance but has near-zero
spatial autocorrelation, unlike real CT/MRI tissue (measured
$\approx 0.6$ autocorrelation at native resolution, versus $\approx 0.0$
for i.i.d.\ paint). An optional multi-octave variant instead sums several
coarse-to-fine random fields before painting, producing spatially
correlated texture that better matches real tissue statistics while
leaving the per-class intensity variance unchanged.
